\chapter{PHÂN TÍCH ĐỘ NHẠY CỤC BỘ}
\label{ch:sensitivity}

\textit{(Nội dung chi tiết của chương này sẽ được phát triển đầy đủ trong giai đoạn tiếp theo):}

\section{Phương pháp vi phân trực tiếp (Direct Differential Method)}
\begin{itemize}
    \item Thiết lập hệ phương trình vi phân mở rộng gồm 2 phương trình trạng thái và các phương trình cho véctơ độ nhạy:
    \begin{equation}
    \frac{d}{dt}\left(\frac{\partial u_i}{\partial p_j}\right) = J \cdot S_j + F_j,
    \end{equation}
    với véctơ 8 tham số $p = (r, a, b, E_1, e, c, d, E_2)$.
    \item Tính ma trận Jacobian $J(x, y)$ và ma trận đạo hàm riêng của vế phải $F$.
\end{itemize}

\section{Chỉ số độ nhạy tương đối và xếp hạng tham số}
\begin{itemize}
    \item Tính chuẩn Euclidean của ma trận độ nhạy tương đối:
    \begin{equation}
    \|\bar{S}_{ij}\|_2 = \sqrt{\sum_{k=1}^N \left(\frac{\partial u_i}{\partial p_j}(t_k) \frac{p_j}{u_i(t_k)}\right)^2}.
    \end{equation}
    \item \textbf{Đối với cá vược sọc:} Nhạy cảm nhất với tốc độ tăng trưởng nội tại của cá mòi dầu $r$ ($\|\cdot\|_2 \approx 1.0080$), hệ số săn bắt $b$ ($\approx 1.0004$) và nỗ lực đánh bắt cá mòi dầu $E_1$ ($1.9647 E_1/q_1$).
    \item \textbf{Đối với cá mòi dầu:} Nhạy cảm nhất với hiệu suất chuyển hóa dinh dưỡng $d$ ($\approx 0.9985$), tỷ lệ tử vong tự nhiên của kẻ săn mồi $e$ ($\approx 0.8886$) và nỗ lực đánh bắt cá vược sọc $E_2$ ($0.7328 E_2/q_2$).
    \item Thảo luận ý nghĩa thực nghiệm đối với các nhà quản lý sinh thái.
\end{itemize}
